% --- Archon physics formalization source begin ---
% archon:physics
% archon:covers PhyXMiniProblems/problem_phyx_mini_0623.lean
% archon:source-report reports/phyx_mini/problem_phyx_mini_0623.source.json

\chapter{Physics problem phyx\_mini\_0623}
\label{ch:PhyXMiniProblems_problem_phyx_mini_0623}

\paragraph{Problem source.}
\#\# Physical scenario

The wave nature of electrons is manifested in experiments where an electron beam interacts with the atoms on the surface of a solid, especially crystals. By studying the angular distribution of the diffracted electrons, one can indirectly measure the geometrical arrangement of atoms. Assume that the electrons strike perpendicular to the surface of a solid (see figure), and that their energy is low, \textbackslash{}(\textbackslash{}mathrm\{KE\}=100\textbackslash{},\textbackslash{}mathrm\{eV\}\textbackslash{}), so that they interact only with the surface layer of atoms. The smallest angle at which a diffraction maximum occurs is at \textbackslash{}(24\textasciicircum{}\textbackslash{}circ\textbackslash{}).

\#\# Question

What is the separation \textbackslash{}(d\textbackslash{}) between the atoms on the surface?

\#\# Answer choices

- A: A: 0.20nm

- B: B: 0.25nm

- C: C: 0.30nm

- D: D: 0.35nm

\#\# Figure caption (auxiliary; use the image as primary evidence)

The image is a diagram illustrating the diffraction of an electron beam off of a crystalline structure. Key elements include:

1. **Dots in a grid pattern**: Represent atoms in a crystalline lattice.
2. **Incident electron beam**: Shown as a vertical arrow directed towards the lattice.
3. **Scattered beams**: Two arrows showing paths of electrons being deflected at an angle θ from the surface.
4. **Angle θ**: The angle between the incident beam and the scattered beams.
5. **Labeled distances**: 
   - \textbackslash{}( d \textbackslash{}): The distance between two adjacent atoms.
   - \textbackslash{}( d \textbackslash{}sin \textbackslash{}theta \textbackslash{}): Shown as the horizontal component of the path of one of the scattered beams.
6. **Text**:
   - "Incident electron beam" next to the incoming electron beam.
   - "θ" labeling the angles between beams and lattice lines.

The diagram illustrates the concept of electron diffraction and the relationship between the incident angle, the lattice spacing, and the path difference.

\#\# Dataset metadata

Quantum Phenomena; Physical Model Grounding Reasoning, Multi-Formula Reasoning

\paragraph{Recorded answer/context.}
C: C: 0.30nm

\paragraph{Figure/image path.}
/root/proposal\_for\_physic/science-mango/phyx\_mini\_run/phyx\_data/test\_image/623.png

\paragraph{Formalization target.}
create a compiling Lean file with sorry bodies at `PhyXMiniProblems/problem\_phyx\_mini\_0623.lean`. The Lean declarations must preserve the physical quantities, dimensions or dimensional roles, figure labels, governing-law hypotheses, and final relation expressed by this problem.
Use Mathlib/Physlib names found through LeanExplore where available. If a physics API is missing, introduce faithful local abstractions rather than scalar placeholder aliases.

\begin{theorem}[Physics formalization target]
\label{thm:physics:phyx_mini_0623:target}
\lean{PhyXMiniProblems.ProblemPhyXMini0623.problem_phyx_mini_0623}
\uses{def:physics:phyx-mini-0623:phyxminiproblems-problemphyxmini0623-lengthinmeters, def:physics:phyx-mini-0623:phyxminiproblems-problemphyxmini0623-lengthinnanometers, def:physics:phyx-mini-0623:phyxminiproblems-problemphyxmini0623-energyinjoules, def:physics:phyx-mini-0623:phyxminiproblems-problemphyxmini0623-massinkilograms, def:physics:phyx-mini-0623:phyxminiproblems-problemphyxmini0623-actioninjouleseconds, def:physics:phyx-mini-0623:phyxminiproblems-problemphyxmini0623-speedoflightinmeterspersecond, def:physics:phyx-mini-0623:phyxminiproblems-problemphyxmini0623-degreestoradians, def:physics:phyx-mini-0623:phyxminiproblems-problemphyxmini0623-electronsurfacediffractionsetup, def:physics:phyx-mini-0623:phyxminiproblems-problemphyxmini0623-matcheselectronsurfacediffractionscenario, def:physics:phyx-mini-0623:phyxminiproblems-problemphyxmini0623-matchesproblemreadouts, def:physics:phyx-mini-0623:phyxminiproblems-problemphyxmini0623-matchessuppliedsurfacediffractionfigure, def:physics:phyx-mini-0623:phyxminiproblems-problemphyxmini0623-usesstandardelectronreferencedata, def:physics:phyx-mini-0623:phyxminiproblems-problemphyxmini0623-hasphysicalsurfacediffractionparameters, def:physics:phyx-mini-0623:phyxminiproblems-problemphyxmini0623-satisfieselectronsurfacediffractionlaws, def:physics:phyx-mini-0623:phyxminiproblems-problemphyxmini0623-displayedseparationinnanometers, def:physics:phyx-mini-0623:phyxminiproblems-problemphyxmini0623-recordeddatasetanswer, def:physics:phyx-mini-0623:phyxminiproblems-problemphyxmini0623-roundstohundredthnanometer, def:physics:phyx-mini-0623:phyxminiproblems-problemphyxmini0623-isuniquematchingdisplayedseparation}
The assigned autoformalize agent should translate this physics problem into Lean declarations in the covered file, with theorem and lemma proof bodies written as `by sorry`.
\end{theorem}
\begin{proof}
This is an autoformalization task, not a proof task. Produce statements that can later be proved by the physics prover without weakening the physical model.
\end{proof}

% --- Archon declaration topology begin ---
\section{Declaration topology}
\begin{definition}[Length Quantity]
  \label{def:physics:phyx-mini-0623:phyxminiproblems-problemphyxmini0623-lengthquantity}
  \lean{PhyXMiniProblems.ProblemPhyXMini0623.LengthQuantity}
  A nonnegative, unit-independent physical length
\end{definition}
\begin{proof}
  This is the declared model definition of Length Quantity.
\end{proof}
\begin{definition}[Mass Quantity]
  \label{def:physics:phyx-mini-0623:phyxminiproblems-problemphyxmini0623-massquantity}
  \lean{PhyXMiniProblems.ProblemPhyXMini0623.MassQuantity}
  A nonnegative, unit-independent physical mass
\end{definition}
\begin{proof}
  This is the declared model definition of Mass Quantity.
\end{proof}
\begin{definition}[momentum Dimension]
  \label{def:physics:phyx-mini-0623:phyxminiproblems-problemphyxmini0623-momentumdimension}
  \lean{PhyXMiniProblems.ProblemPhyXMini0623.momentumDimension}
  The dimension mass * length / time of a momentum magnitude
\end{definition}
\begin{proof}
  This is the declared model definition of momentum Dimension.
\end{proof}
\begin{definition}[Momentum Magnitude Quantity]
  \label{def:physics:phyx-mini-0623:phyxminiproblems-problemphyxmini0623-momentummagnitudequantity}
  \lean{PhyXMiniProblems.ProblemPhyXMini0623.MomentumMagnitudeQuantity}
  \uses{def:physics:phyx-mini-0623:phyxminiproblems-problemphyxmini0623-momentumdimension}
  A nonnegative, unit-independent physical momentum magnitude
\end{definition}
\begin{proof}
  This is the declared model definition of Momentum Magnitude Quantity.
\end{proof}
\begin{definition}[action Dimension]
  \label{def:physics:phyx-mini-0623:phyxminiproblems-problemphyxmini0623-actiondimension}
  \lean{PhyXMiniProblems.ProblemPhyXMini0623.actionDimension}
  The dimension mass * length² / time of an action
\end{definition}
\begin{proof}
  This is the declared model definition of action Dimension.
\end{proof}
\begin{definition}[Action Quantity]
  \label{def:physics:phyx-mini-0623:phyxminiproblems-problemphyxmini0623-actionquantity}
  \lean{PhyXMiniProblems.ProblemPhyXMini0623.ActionQuantity}
  \uses{def:physics:phyx-mini-0623:phyxminiproblems-problemphyxmini0623-actiondimension}
  A nonnegative, unit-independent physical action
\end{definition}
\begin{proof}
  This is the declared model definition of Action Quantity.
\end{proof}
\begin{definition}[length Readout]
  \label{def:physics:phyx-mini-0623:phyxminiproblems-problemphyxmini0623-lengthreadout}
  \lean{PhyXMiniProblems.ProblemPhyXMini0623.lengthReadout}
  \uses{def:physics:phyx-mini-0623:phyxminiproblems-problemphyxmini0623-lengthquantity}
  Read a physical length as a real scalar in a selected length unit
\end{definition}
\begin{proof}
  This is the declared model definition of length Readout.
\end{proof}
\begin{definition}[length In Meters]
  \label{def:physics:phyx-mini-0623:phyxminiproblems-problemphyxmini0623-lengthinmeters}
  \lean{PhyXMiniProblems.ProblemPhyXMini0623.lengthInMeters}
  \uses{def:physics:phyx-mini-0623:phyxminiproblems-problemphyxmini0623-lengthquantity, def:physics:phyx-mini-0623:phyxminiproblems-problemphyxmini0623-lengthreadout}
  Metre readout of a physical length
\end{definition}
\begin{proof}
  This is the declared model definition of length In Meters.
\end{proof}
\begin{definition}[length In Nanometers]
  \label{def:physics:phyx-mini-0623:phyxminiproblems-problemphyxmini0623-lengthinnanometers}
  \lean{PhyXMiniProblems.ProblemPhyXMini0623.lengthInNanometers}
  \uses{def:physics:phyx-mini-0623:phyxminiproblems-problemphyxmini0623-lengthquantity, def:physics:phyx-mini-0623:phyxminiproblems-problemphyxmini0623-lengthreadout}
  Nanometre readout used for the four displayed separations
\end{definition}
\begin{proof}
  This is the declared model definition of length In Nanometers.
\end{proof}
\begin{definition}[energy In Joules]
  \label{def:physics:phyx-mini-0623:phyxminiproblems-problemphyxmini0623-energyinjoules}
  \lean{PhyXMiniProblems.ProblemPhyXMini0623.energyInJoules}
  Coherent-SI energy readout in joules
\end{definition}
\begin{proof}
  This is the declared model definition of energy In Joules.
\end{proof}
\begin{definition}[energy In Electron Volts]
  \label{def:physics:phyx-mini-0623:phyxminiproblems-problemphyxmini0623-energyinelectronvolts}
  \lean{PhyXMiniProblems.ProblemPhyXMini0623.energyInElectronVolts}
  \uses{def:physics:phyx-mini-0623:phyxminiproblems-problemphyxmini0623-energyinjoules}
  Electron-volt readout grounded by Physlib's dimensionful calibrated DimEnergy.electronVolt, rather than by treating energy as a bare scalar
\end{definition}
\begin{proof}
  This is the declared model definition of energy In Electron Volts.
\end{proof}
\begin{definition}[mass In Kilograms]
  \label{def:physics:phyx-mini-0623:phyxminiproblems-problemphyxmini0623-massinkilograms}
  \lean{PhyXMiniProblems.ProblemPhyXMini0623.massInKilograms}
  \uses{def:physics:phyx-mini-0623:phyxminiproblems-problemphyxmini0623-massquantity}
  Coherent-SI mass readout in kilograms
\end{definition}
\begin{proof}
  This is the declared model definition of mass In Kilograms.
\end{proof}
\begin{definition}[momentum In Kilogram Meters Per Second]
  \label{def:physics:phyx-mini-0623:phyxminiproblems-problemphyxmini0623-momentuminkilogrammeterspersecond}
  \lean{PhyXMiniProblems.ProblemPhyXMini0623.momentumInKilogramMetersPerSecond}
  \uses{def:physics:phyx-mini-0623:phyxminiproblems-problemphyxmini0623-momentummagnitudequantity}
  Coherent-SI momentum-magnitude readout in kilogram-metres per second
\end{definition}
\begin{proof}
  This is the declared model definition of momentum In Kilogram Meters Per Second.
\end{proof}
\begin{definition}[action In Joule Seconds]
  \label{def:physics:phyx-mini-0623:phyxminiproblems-problemphyxmini0623-actioninjouleseconds}
  \lean{PhyXMiniProblems.ProblemPhyXMini0623.actionInJouleSeconds}
  \uses{def:physics:phyx-mini-0623:phyxminiproblems-problemphyxmini0623-actionquantity}
  Coherent-SI action readout in joule-seconds
\end{definition}
\begin{proof}
  This is the declared model definition of action In Joule Seconds.
\end{proof}
\begin{definition}[speed Of Light In Meters Per Second]
  \label{def:physics:phyx-mini-0623:phyxminiproblems-problemphyxmini0623-speedoflightinmeterspersecond}
  \lean{PhyXMiniProblems.ProblemPhyXMini0623.speedOfLightInMetersPerSecond}
  Coherent-SI readout of Physlib's exact vacuum light-speed constant
\end{definition}
\begin{proof}
  This is the declared model definition of speed Of Light In Meters Per Second.
\end{proof}
\begin{definition}[degrees To Radians]
  \label{def:physics:phyx-mini-0623:phyxminiproblems-problemphyxmini0623-degreestoradians}
  \lean{PhyXMiniProblems.ProblemPhyXMini0623.degreesToRadians}
  Convert the displayed degree measure of an angle to radians
\end{definition}
\begin{proof}
  This is the declared model definition of degrees To Radians.
\end{proof}
\begin{definition}[Particle Species]
  \label{def:physics:phyx-mini-0623:phyxminiproblems-problemphyxmini0623-particlespecies}
  \lean{PhyXMiniProblems.ProblemPhyXMini0623.ParticleSpecies}
  Particle species in the incident beam
\end{definition}
\begin{proof}
  This is the declared model definition of Particle Species.
\end{proof}
\begin{definition}[Surface Kind]
  \label{def:physics:phyx-mini-0623:phyxminiproblems-problemphyxmini0623-surfacekind}
  \lean{PhyXMiniProblems.ProblemPhyXMini0623.SurfaceKind}
  Kind of solid surface struck by the beam
\end{definition}
\begin{proof}
  This is the declared model definition of Surface Kind.
\end{proof}
\begin{definition}[Incidence Geometry]
  \label{def:physics:phyx-mini-0623:phyxminiproblems-problemphyxmini0623-incidencegeometry}
  \lean{PhyXMiniProblems.ProblemPhyXMini0623.IncidenceGeometry}
  Direction of incidence relative to the surface
\end{definition}
\begin{proof}
  This is the declared model definition of Incidence Geometry.
\end{proof}
\begin{definition}[Figure Beam]
  \label{def:physics:phyx-mini-0623:phyxminiproblems-problemphyxmini0623-figurebeam}
  \lean{PhyXMiniProblems.ProblemPhyXMini0623.FigureBeam}
  Beam arrows represented in the primary diagram
\end{definition}
\begin{proof}
  This is the declared model definition of Figure Beam.
\end{proof}
\begin{definition}[Figure Length Label]
  \label{def:physics:phyx-mini-0623:phyxminiproblems-problemphyxmini0623-figurelengthlabel}
  \lean{PhyXMiniProblems.ProblemPhyXMini0623.FigureLengthLabel}
  Length labels printed in the primary diagram
\end{definition}
\begin{proof}
  This is the declared model definition of Figure Length Label.
\end{proof}
\begin{definition}[Surface Diffraction Figure]
  \label{def:physics:phyx-mini-0623:phyxminiproblems-problemphyxmini0623-surfacediffractionfigure}
  \lean{PhyXMiniProblems.ProblemPhyXMini0623.SurfaceDiffractionFigure}
  \uses{def:physics:phyx-mini-0623:phyxminiproblems-problemphyxmini0623-figurebeam, def:physics:phyx-mini-0623:phyxminiproblems-problemphyxmini0623-figurelengthlabel}
  Qualitative information transcribed from image 623.png. The diagram is schematic: it supplies labels and incidence/scattering geometry, but no additional numerical length measurement
\end{definition}
\begin{proof}
  This is the declared model definition of Surface Diffraction Figure.
\end{proof}
\begin{definition}[Electron Surface Diffraction Setup]
  \label{def:physics:phyx-mini-0623:phyxminiproblems-problemphyxmini0623-electronsurfacediffractionsetup}
  \lean{PhyXMiniProblems.ProblemPhyXMini0623.ElectronSurfaceDiffractionSetup}
  \uses{def:physics:phyx-mini-0623:phyxminiproblems-problemphyxmini0623-lengthquantity, def:physics:phyx-mini-0623:phyxminiproblems-problemphyxmini0623-massquantity, def:physics:phyx-mini-0623:phyxminiproblems-problemphyxmini0623-momentummagnitudequantity, def:physics:phyx-mini-0623:phyxminiproblems-problemphyxmini0623-actionquantity, def:physics:phyx-mini-0623:phyxminiproblems-problemphyxmini0623-particlespecies, def:physics:phyx-mini-0623:phyxminiproblems-problemphyxmini0623-surfacekind, def:physics:phyx-mini-0623:phyxminiproblems-problemphyxmini0623-incidencegeometry, def:physics:phyx-mini-0623:phyxminiproblems-problemphyxmini0623-surfacediffractionfigure}
  Independent physical quantities in the experiment. pathDifferenceAt is the outgoing-ray path difference for an angle measured in degrees, and isDiffractionMaximum order angle records an observed maximum without defining either the wavelength or atom separation
\end{definition}
\begin{proof}
  This is the declared model definition of Electron Surface Diffraction Setup.
\end{proof}
\begin{definition}[Matches Electron Surface Diffraction Scenario]
  \label{def:physics:phyx-mini-0623:phyxminiproblems-problemphyxmini0623-matcheselectronsurfacediffractionscenario}
  \lean{PhyXMiniProblems.ProblemPhyXMini0623.MatchesElectronSurfaceDiffractionScenario}
  \uses{def:physics:phyx-mini-0623:phyxminiproblems-problemphyxmini0623-electronsurfacediffractionsetup}
  Qualitative physical scenario stated in the problem
\end{definition}
\begin{proof}
  This is the declared model definition of Matches Electron Surface Diffraction Scenario.
\end{proof}
\begin{definition}[Matches Problem Readouts]
  \label{def:physics:phyx-mini-0623:phyxminiproblems-problemphyxmini0623-matchesproblemreadouts}
  \lean{PhyXMiniProblems.ProblemPhyXMini0623.MatchesProblemReadouts}
  \uses{def:physics:phyx-mini-0623:phyxminiproblems-problemphyxmini0623-energyinelectronvolts, def:physics:phyx-mini-0623:phyxminiproblems-problemphyxmini0623-electronsurfacediffractionsetup}
  The two numerical readouts stated in the prose. The designation of the smallest positive maximum as order one is the usual diffraction-order naming for this normal-incidence setup. No wavelength or atom-separation value occurs here
\end{definition}
\begin{proof}
  This is the declared model definition of Matches Problem Readouts.
\end{proof}
\begin{definition}[Matches Supplied Surface Diffraction Figure]
  \label{def:physics:phyx-mini-0623:phyxminiproblems-problemphyxmini0623-matchessuppliedsurfacediffractionfigure}
  \lean{PhyXMiniProblems.ProblemPhyXMini0623.MatchesSuppliedSurfaceDiffractionFigure}
  \uses{def:physics:phyx-mini-0623:phyxminiproblems-problemphyxmini0623-surfacediffractionfigure}
  Facts directly visible in the supplied surface-diffraction diagram
\end{definition}
\begin{proof}
  This is the declared model definition of Matches Supplied Surface Diffraction Figure.
\end{proof}
\begin{definition}[Uses Standard Electron Reference Data]
  \label{def:physics:phyx-mini-0623:phyxminiproblems-problemphyxmini0623-usesstandardelectronreferencedata}
  \lean{PhyXMiniProblems.ProblemPhyXMini0623.UsesStandardElectronReferenceData}
  \uses{def:physics:phyx-mini-0623:phyxminiproblems-problemphyxmini0623-massinkilograms, def:physics:phyx-mini-0623:phyxminiproblems-problemphyxmini0623-actioninjouleseconds, def:physics:phyx-mini-0623:phyxminiproblems-problemphyxmini0623-electronsurfacediffractionsetup}
  Standard reference data needed by the exact relativistic de Broglie calculation. Physlib provides Constants.ℏ in joule-seconds, so the ordinary Planck action has readout 2πℏ. Physlib has no electron-mass constant, hence the CODATA electron mass is a calibrated premise. The exact light-speed constant is grounded by DimSpeed.speedOfLight. None of these reference values contains the requested atom separation
\end{definition}
\begin{proof}
  This is the declared model definition of Uses Standard Electron Reference Data.
\end{proof}
\begin{definition}[Has Physical Surface Diffraction Parameters]
  \label{def:physics:phyx-mini-0623:phyxminiproblems-problemphyxmini0623-hasphysicalsurfacediffractionparameters}
  \lean{PhyXMiniProblems.ProblemPhyXMini0623.HasPhysicalSurfaceDiffractionParameters}
  \uses{def:physics:phyx-mini-0623:phyxminiproblems-problemphyxmini0623-lengthinmeters, def:physics:phyx-mini-0623:phyxminiproblems-problemphyxmini0623-energyinjoules, def:physics:phyx-mini-0623:phyxminiproblems-problemphyxmini0623-massinkilograms, def:physics:phyx-mini-0623:phyxminiproblems-problemphyxmini0623-momentuminkilogrammeterspersecond, def:physics:phyx-mini-0623:phyxminiproblems-problemphyxmini0623-actioninjouleseconds, def:physics:phyx-mini-0623:phyxminiproblems-problemphyxmini0623-speedoflightinmeterspersecond, def:physics:phyx-mini-0623:phyxminiproblems-problemphyxmini0623-electronsurfacediffractionsetup}
  Positivity and angle-range conditions selecting the physical regime
\end{definition}
\begin{proof}
  This is the declared model definition of Has Physical Surface Diffraction Parameters.
\end{proof}
\begin{definition}[Satisfies Electron Surface Diffraction Laws]
  \label{def:physics:phyx-mini-0623:phyxminiproblems-problemphyxmini0623-satisfieselectronsurfacediffractionlaws}
  \lean{PhyXMiniProblems.ProblemPhyXMini0623.SatisfiesElectronSurfaceDiffractionLaws}
  \uses{def:physics:phyx-mini-0623:phyxminiproblems-problemphyxmini0623-lengthinmeters, def:physics:phyx-mini-0623:phyxminiproblems-problemphyxmini0623-energyinjoules, def:physics:phyx-mini-0623:phyxminiproblems-problemphyxmini0623-massinkilograms, def:physics:phyx-mini-0623:phyxminiproblems-problemphyxmini0623-momentuminkilogrammeterspersecond, def:physics:phyx-mini-0623:phyxminiproblems-problemphyxmini0623-actioninjouleseconds, def:physics:phyx-mini-0623:phyxminiproblems-problemphyxmini0623-speedoflightinmeterspersecond, def:physics:phyx-mini-0623:phyxminiproblems-problemphyxmini0623-degreestoradians, def:physics:phyx-mini-0623:phyxminiproblems-problemphyxmini0623-electronsurfacediffractionsetup}
  The governing relations used by the textbook calculation: exact relativistic energy--momentum dispersion gives (K + m c²)² = (p c)² + (m c²)²; de Broglie's law gives λ p = h; the figure geometry gives the path difference d sin θ; at an order-n maximum, that path difference is n λ. The last two equations hold for every angle/order supplied to them. They do not specialize the atom separation, wavelength, or answer label to this problem's requested result.
\end{definition}
\begin{proof}
  This is the declared model definition of Satisfies Electron Surface Diffraction Laws.
\end{proof}
\begin{definition}[Answer Choice]
  \label{def:physics:phyx-mini-0623:phyxminiproblems-problemphyxmini0623-answerchoice}
  \lean{PhyXMiniProblems.ProblemPhyXMini0623.AnswerChoice}
  Labels attached to the four displayed atom-separation choices
\end{definition}
\begin{proof}
  This is the declared model definition of Answer Choice.
\end{proof}
\begin{definition}[displayed Separation In Nanometers]
  \label{def:physics:phyx-mini-0623:phyxminiproblems-problemphyxmini0623-displayedseparationinnanometers}
  \lean{PhyXMiniProblems.ProblemPhyXMini0623.displayedSeparationInNanometers}
  \uses{def:physics:phyx-mini-0623:phyxminiproblems-problemphyxmini0623-answerchoice}
  Separation in nanometres printed beside each answer label
\end{definition}
\begin{proof}
  This is the declared model definition of displayed Separation In Nanometers.
\end{proof}
\begin{definition}[recorded Dataset Answer]
  \label{def:physics:phyx-mini-0623:phyxminiproblems-problemphyxmini0623-recordeddatasetanswer}
  \lean{PhyXMiniProblems.ProblemPhyXMini0623.recordedDatasetAnswer}
  \uses{def:physics:phyx-mini-0623:phyxminiproblems-problemphyxmini0623-answerchoice}
  The answer label recorded by the supplied dataset
\end{definition}
\begin{proof}
  This is the declared model definition of recorded Dataset Answer.
\end{proof}
\begin{definition}[Rounds To Hundredth Nanometer]
  \label{def:physics:phyx-mini-0623:phyxminiproblems-problemphyxmini0623-roundstohundredthnanometer}
  \lean{PhyXMiniProblems.ProblemPhyXMini0623.RoundsToHundredthNanometer}
  Agreement with a separation printed to two decimal places in nanometres
\end{definition}
\begin{proof}
  This is the declared model definition of Rounds To Hundredth Nanometer.
\end{proof}
\begin{definition}[Matches Displayed Separation]
  \label{def:physics:phyx-mini-0623:phyxminiproblems-problemphyxmini0623-matchesdisplayedseparation}
  \lean{PhyXMiniProblems.ProblemPhyXMini0623.MatchesDisplayedSeparation}
  \uses{def:physics:phyx-mini-0623:phyxminiproblems-problemphyxmini0623-lengthinnanometers, def:physics:phyx-mini-0623:phyxminiproblems-problemphyxmini0623-electronsurfacediffractionsetup, def:physics:phyx-mini-0623:phyxminiproblems-problemphyxmini0623-answerchoice, def:physics:phyx-mini-0623:phyxminiproblems-problemphyxmini0623-displayedseparationinnanometers, def:physics:phyx-mini-0623:phyxminiproblems-problemphyxmini0623-roundstohundredthnanometer}
  A displayed choice agrees with the atom separation to the shown precision
\end{definition}
\begin{proof}
  This is the declared model definition of Matches Displayed Separation.
\end{proof}
\begin{definition}[Is Unique Matching Displayed Separation]
  \label{def:physics:phyx-mini-0623:phyxminiproblems-problemphyxmini0623-isuniquematchingdisplayedseparation}
  \lean{PhyXMiniProblems.ProblemPhyXMini0623.IsUniqueMatchingDisplayedSeparation}
  \uses{def:physics:phyx-mini-0623:phyxminiproblems-problemphyxmini0623-electronsurfacediffractionsetup, def:physics:phyx-mini-0623:phyxminiproblems-problemphyxmini0623-answerchoice, def:physics:phyx-mini-0623:phyxminiproblems-problemphyxmini0623-matchesdisplayedseparation}
  A choice is the unique displayed separation matching the modeled result
\end{definition}
\begin{proof}
  This is the declared model definition of Is Unique Matching Displayed Separation.
\end{proof}
\begin{lemma}[matter Wavelength from relativistic Energy]
  \label{lem:physics:phyx-mini-0623:phyxminiproblems-problemphyxmini0623-matterwavelength-from-relativisticenergy}
  \lean{PhyXMiniProblems.ProblemPhyXMini0623.matterWavelength_from_relativisticEnergy}
  \uses{def:physics:phyx-mini-0623:phyxminiproblems-problemphyxmini0623-lengthinmeters, def:physics:phyx-mini-0623:phyxminiproblems-problemphyxmini0623-energyinjoules, def:physics:phyx-mini-0623:phyxminiproblems-problemphyxmini0623-massinkilograms, def:physics:phyx-mini-0623:phyxminiproblems-problemphyxmini0623-actioninjouleseconds, def:physics:phyx-mini-0623:phyxminiproblems-problemphyxmini0623-speedoflightinmeterspersecond, def:physics:phyx-mini-0623:phyxminiproblems-problemphyxmini0623-electronsurfacediffractionsetup, def:physics:phyx-mini-0623:phyxminiproblems-problemphyxmini0623-hasphysicalsurfacediffractionparameters, def:physics:phyx-mini-0623:phyxminiproblems-problemphyxmini0623-satisfieselectronsurfacediffractionlaws}
  The exact relativistic dispersion and de Broglie laws determine the electron matter wavelength without using any diffraction-spacing conclusion. The term (K / c)² is retained, so this is not the globalized nonrelativistic approximation rejected by the formalization review gate
\end{lemma}
\begin{proof}
  The conclusion follows from the governing relations and figure readouts named in the statement of matter Wavelength from relativistic Energy.
\end{proof}
\begin{lemma}[atom Separation from first Maximum]
  \label{lem:physics:phyx-mini-0623:phyxminiproblems-problemphyxmini0623-atomseparation-from-firstmaximum}
  \lean{PhyXMiniProblems.ProblemPhyXMini0623.atomSeparation_from_firstMaximum}
  \uses{def:physics:phyx-mini-0623:phyxminiproblems-problemphyxmini0623-lengthinnanometers, def:physics:phyx-mini-0623:phyxminiproblems-problemphyxmini0623-degreestoradians, def:physics:phyx-mini-0623:phyxminiproblems-problemphyxmini0623-electronsurfacediffractionsetup, def:physics:phyx-mini-0623:phyxminiproblems-problemphyxmini0623-matchesproblemreadouts, def:physics:phyx-mini-0623:phyxminiproblems-problemphyxmini0623-hasphysicalsurfacediffractionparameters, def:physics:phyx-mini-0623:phyxminiproblems-problemphyxmini0623-satisfieselectronsurfacediffractionlaws}
  At the observed first-order maximum, d sin θ = λ; solving this general relation for d gives the exact spacing formula used before rounding
\end{lemma}
\begin{proof}
  The conclusion follows from the governing relations and figure readouts named in the statement of atom Separation from first Maximum.
\end{proof}
% --- Archon declaration topology end ---
% --- Archon physics formalization source end ---
